Small and medium-sized enterprises are central to employment and structural transformation in Rwanda, yet many operate under market disruption, uneven digitalization, and constrained adaptive capacity. This paper audits and predicts a constructed firm-level fragility classification using the 2023 Rwanda Enterprise Survey. We recover the exact deterministic rule behind a legacy binary fragility indicator and show that it classifies 62 of 358 firms, or 17.3 percent of observations, as fragile. The recovered rule combines product/service innovation, process innovation, website status, and permanent full-time employment, exposing a central construct-validity problem: earlier project drafts label some of the same variables as revenue volatility and digital-platform difficulty. We therefore treat the official Stata labels as authoritative and interpret the outcome as a screening index rather than an observed survival, closure, or financial-distress measure. A least absolute shrinkage and selection operator (LASSO) logistic-regression pipeline is estimated with a held-out test split. The full model achieves an area under the receiver operating characteristic curve (ROC-AUC) of 0.978, but it includes variables used to construct the outcome and is therefore interpreted as screening-rule replication. The leakage-reduced benchmark excludes outcome-rule variables and yields ROC-AUC of 0.770, area under the precision-recall curve (PR-AUC) of 0.428, and accuracy of 61.1 percent. The central contribution is a reproducible audit workflow showing how constructed policy-screening indices should be validated before they are interpreted as measures of firm distress. The paper clarifies the validation, governance, and external-outcome evidence required before SME fragility scores can inform policy targeting.
